\subsection{\DVRFOne, A Concrete Verifiable Pseudorandom Secret Sharing Scheme}
\label{dvrf-concrete}

We now define a concrete VPSS that we call \DVRFOne,
that builds upon the pseudorandom secret sharing scheme as defined by Cramer Damg{\aa}rd, and Ishai~\cite{CramerDI05}.
However, \DVRFOne additionally defines a verify algorithm that is publicly verifiable (assuming the authenticity of the inputs),
as well as an algorithm to combine participant commitments.

As a reminder,
the pseudorandom secret sharing scheme by Cramer, Damg{\aa}rd, and Ishai~\cite{CramerDI05}  itself builds on replicated secret sharing~\cite{MitsuruSN89}.
However, Cramer et al.\ define a mechanism to non-interactively convert shares of a replicated secret sharing scheme to shares of a Shamir-secret-shared value.
They then show how this mechanism can be used as a distributed pseudorandom function;
i.e., given a replicated secret sharing of a random value,
participants can generate Shamir secret shares of an unbounded number of pseudorandom values,
without interaction.

We extend this pseudorandom secret sharing scheme by Cramer et al. and additionally define a public verifiability mechanism.
More specifically,
we define a $\verify$ function that uses the set of participants' commitments to ensure all participants performed the evaluation step  in a consistent manner.
Such consistency checks have been used in prior literature~\cite{BenhamoudaBGHIN21},
but our check is performed over commitments to secret shares,
as opposed to verifying the shares directly.
%
We prove that \DVRFOne is secure assuming a cryptographically secure hash function,
when at minimum $2t-1$ parties participate in evaluation and up to $t-1$ participants are corrupted (honest majority).

We now define \DVRFOne in more detail. We show $\keygen$ for \DVRFOne as a centralized operation,
but this operation can easily be decentralized.
\DVRFOne is additionally defined with respect to a hash function $\Hash$,
where $\Hash: \Zq \times \zo^* \rightarrow \Zq$ is a cryptographically secure hash function.

\begin{itemize}
  \item $\shinem.\setup(\usecp)$: Accepts as input the security parameter $\secp$,
    outputs public parameters $\pp = (\GG, q, g)$,
    where $(\GG, q, g)$ is generated by $\GroupGen$.
    $\pp$ is given as an implicit input to all other algorithms.
    %\ian{Minor nit (I've fixed): you can't say $\shinem[\Hash].\setup$ because $\setup$ outputs $q$ and the definition of $\Hash$ requires $q$ to have already been chosen. So technically, $\Hash$ can only be chosen \emph{after} $\setup$ completes.  Would it be cleaner to just have $\Hash$ be an \emph{output} of $\setup$, and so in $\pp$, and then you don't have to write it explicitly as ``$\shinem[\Hash]$'' anywhere?  Do you rely in the proofs on reprogramming or otherwise messing with the $\Hash$ that's used in $\shinem$?}
    %\chelsea{I set $\Hash$} explicitly because it is defined by the outer
    %function
  \item $\shinem[\Hash].\keygen(n, t, \minparticipants) \rightarrow \bot/(\SK_1, \ldots, \SK_n)$:
    On input the total number of participants $n$,
    the corruption threshold $t$,
    and minimum number of participants $\minparticipants$,
    the dealer performs replicated secret sharing of a random secret $\SK \in \Zq$~\cite{CramerDI05},
    following the following steps:
    \begin{enumerate}
      \item First, the dealer checks if $\minparticipants \geq 2t-1$;
        if the check fails, output $\bot$.
      \item Let $A$ be the set such that $A = {\set{n} \choose t-1}$,
        and let $\gamma$ be the size of $A$;
        i.e, $\gamma = \lvert A \rvert = {n \choose t-1}$.
      \item Generate $\gamma$ secret shares $\{\share_i\}_{i \in \set{\gamma}}$,
        by sampling $\share_i \rgets \Zq, i \in \set{\gamma}$.
        Implicitly,
        the secret $\SK$ is such that $\SK = \sum_{i \in \gamma} \share_i$.

      \item Initialize empty sets $\SK_1 = \emptyset, \ldots, \SK_n = \emptyset$.
      \item For each set $a_i \in A, i \in \set{\gamma}$
        and each participant identifier $j \in \set{n} \setminus a_i$,
        append $\SK_j \gets \SK_j \cup \{ (a_i, \share_i) \}$.
      \item Output $(\SK_1, \ldots, \SK_n)$.
        %\chelsea{We need to use consistent notation below with what we describe in preliminaries.}
        Each $\SK_j$ is a set of size $\delta = {n-1 \choose t-1}$,
        consisting of the tuples $(a_i, \share_i)$.
        Each $a_i$ is itself a set,
        such that $a_i \subset \set{n}$ and $\lvert a_i \rvert = t-1$,
        where if $a_i \in \sk_j$, then $j \not\in a_i$.
    \end{enumerate}

  \item $\shinem[\Hash].\generate(\partyid, \SK_\partyid, \vrfin) \rightarrow (\vrfprivout_\partyid, \vrfpubout_\partyid)$:
    On input participant identifier $\partyid \in \set{n}$,
    secret key share $\SK_\partyid$,
    and input $\vrfin \in \{0,1\}^*$,
    perform the following steps:
    \begin{enumerate}
      \item Parse $\{ (a_i, \share_i) \}_{i \in \set{\numsets}} \gets \SK_\partyid$.
      \item For each set $a_i$,
        let $\cramerpoly_{a_i}(x)$ be the degree $t-1$ polynomial defined by the set $a_i$,
        as given in Equation~\ref{eq:cramerpoly}:

        \begin{equation} \label{eq:cramerpoly}
          \cramerpoly_{a_i}(x) = \prod_{j \in a_i}\frac{j - x}{j}
        \end{equation}

        Note that $\cramerpoly_{a_i}(j) = 0$ for all $j \in a_i$ and $\cramerpoly_{a_i}(0) = 1$.

      \item Obtain the share via Equation~\ref{eq:vrf-out-eval}:
    \begin{equation} \label{eq:vrf-out-eval}
      \vrfprivout_\partyid \gets \sum_{i \in \set{\numsets}} \Hash(\share_i, \vrfin) \cdot \cramerpoly_{a_i}(\partyid)
    \end{equation}
    Note that the $\cramerpoly_{a_i}(\partyid)$ values can be precomputed by participant $k$, as they are independent of $\vrfin$, so this step requires $\numsets$ evaluations of $\Hash$ and $\numsets$ multiplications and additions in $\Zq$.

      \item Derive the commitment to $\vrfprivout_\partyid$ as $\vrfpubout_\partyid \gets g^{\vrfprivout_\partyid}$.
      \item Output $(\vrfprivout_\partyid, \vrfpubout_\partyid)$.
    \end{enumerate}

  \item $\shinem[\Hash].\verify(t, \minparticipants, \coalition, \{ \vrfpubout_j \}_{j \in \coalition}) \rightarrow \zo$:
    Perform the following steps:
    \begin{enumerate}
      \item If $\lvert \coalition \rvert \not\geq \minparticipants$,
        return $0$.
      \item Otherwise, define $\polycomm = (\polycomm_0, \polycomm_1, \ldots, \polycomm_{\lvert\coalition\rvert-1})$
        to be the tuple of $\lvert \coalition \rvert$ commitments to \emph{coefficients}
        of a polynomial $f$ defined ``in the exponent,''
        where each $\vrfpubout_i, i \in \coalition$ is a commitment to a \emph{point} on the same polynomial $f$.
      Each $\polycomm_i$ can be derived via Equation~\ref{eq:eval-vrfcomm}:
        \begin{equation} \label{eq:eval-vrfcomm}
          \polycomm_i = \prod_{j \in \coalition} \vrfpubout_j^{L_{j,i}}
        \end{equation}
        where $L_{j,i}$ is the coefficient of  the $i\supth$ term $x^i$ of the $j\supth$ Lagrange polynomial $L_j(x)$ for the coalition $\coalition$,
        as described in Section~\ref{prelim:notation}.

      \item Let $\ell = \lvert \coalition \rvert - t$.
        The verifier now checks that all participants followed the protocol honestly,
        by checking in the exponent that their shares lie on a polynomial of degree $t-1$.
        In particular,
        the verifier ensures that $\polycomm$ is a commitment to a polynomial of degree at most $t-1$,\footnote{When at least $t$ parties follow the protocol honestly, their shares will completely define this polynomial.}
        by checking the last $\ell$ entries in $\polycomm$ are equal to the identity of $\GG$:\footnote{As an optimization, note that $\polycomm_1, \ldots, \polycomm_{t-1}$ need not be computed.}
        \begin{equation} \label{eq:poly-check}
          \polycomm_{t-1 + i} = \identityG, \forall\ i \in \set{\ell}
        \end{equation}
      \item If the check in Equation~\ref{eq:poly-check} holds,
        output $1$,
        otherwise, output $0$.

    \end{enumerate}

  \item $\shinem[\Hash].\combinepub(t, \minparticipants, \coalition, \{ \vrfpubout_j \}_{j \in \coalition}) \rightarrow \vrfpubout $:
    Accepts as input the corruption threshold $t$,
    minimum participants $\minparticipants$,
    the coalition $\coalition$,
    and set of (verified) commitments $\{ \vrfpubout_j \}_{j \in \coalition})$.
    The commitments can be combined as in Equation~\ref{eq:combine-proof},
    where the coefficients $\lambda_j$ are determined by the coalition $\coalition$.
        \begin{equation} \label{eq:combine-proof}
          \vrfpubout \gets \prod_{j \in \coalition} \vrfpubout_j^{\lambda_j} = \polycomm_0
        \end{equation}
    Output $\vrfpubout$.

  \item $\shinem[\Hash].\recover(t, \minparticipants, \coalition, \{ \vrfprivout_j \}_{j \in \coalition}) \rightarrow  \bot/(\vrfprivout, \vrfpubout)$:
    Receives as input the corruption threshold $t$,
    the minimum number of participants $\minparticipants$,
    the coalition $\coalition$,
    and the set of pseudorandom secret shares $\{ \vrfprivout_j \}_{j \in \coalition}$.
    Performs the following steps:
    \begin{enumerate}
      \item If $\lvert \coalition \rvert < \minparticipants$,
        or $\coalition \not\subset \set{n}$,
        output $\bot$.
      \item For each $\vrfprivout_j, j \in \coalition$,
        derive $\vrfpubout_j \gets g^{\vrfpubout_j}$.
      \item If $\DVRFOnem.\verify(t, \minparticipants, \coalition, \{ \vrfpubout_j\}_{j \in \coalition}) \neq 1$,
        then output $\bot$.
      \item Otherwise, the shares can be combined as in Equation~\ref{eq:combine-out},
    where the $\lambda_j$ are determined by the coalition $\coalition$.
        \begin{equation} \label{eq:combine-out}
          \vrfprivout \gets \sum_{j \in \coalition} \vrfprivout_j \cdot \lambda_j
        \end{equation}
      \item Derive the commitment to $\vrfprivout$ as $\vrfpubout \gets g^{\vrfprivout}$.
      \item Output $(\vrfprivout, \vrfpubout)$.
    \end{enumerate}

\end{itemize}

\paragraph{Correctness.}
By correctness of pseudorandom secret sharing as defined by Cramer et al.~\cite{CramerDI05},
when key generation is honestly performed,
the $\generate$ algorithm produces secret shares $\vrfprivout_i = f(i)$
that are points on a degree $t-1$ Shamir secret sharing polynomial $f(x) = \sum_{i \in \set{\gamma}} \Hash(\share_i,\vrfin) \cramerpoly_{a_i}(x)$.
The combined commitment $\vrfprivout = \sum_{i \in \coalition}
\vrfprivout_i \lambda_i$ for $\coalition \subseteq \set{n}, \lvert \coalition
\rvert \geq 2t-1$ is a commitment to a pseudorandom secret value
$\vrfprivout = f(0) = \sum_{i \in \set{\gamma}} \Hash(\share_i,\vrfin)$.
$\combinepub$ and $\recover$ simply perform polynomial interpolation of these values with respect to the set of participants,
and so are likewise correct.

$\verify$ is correct because honest parties output commitments $\vrfpubout_i  = g^{f(i)}$ for the degree $t-1$ polynomial $f$ defined above.
Therefore, if all parties are honest, the coefficients of $x^t$, \dots, $x^{\lvert\coalition\rvert-1}$ of $f(x)$ will be 0, so the commitments $\polycomm_t$, \dots, $\polycomm_{\lvert\coalition\rvert-1}$ to those coefficients will be $\identityG$
(the identity element of $\GG$), and so $\verify$ will output 1.
% CK- for correctness, we assume all parties follow the protocol honestly.
%When parties follow the protocol honestly,
%the polynomial interpolating $(\vrfpubout_j$ values in the exponent will be exactly $f(x)$ if its degree is at most $t-1$.
%(and so all parties output the correct values) and will have degree at least $t$ if any corrupt party output an incorrect value.

\subsubsection{Security}
\label{shine-security-thm}

\DVRFOne is verifiable, unique, and pseudorandom.
We give the corresponding proofs next.

